\section{Curvature reads under one actual Gaussian reader}
\label{sec:actual-reader-guarded-curvature}

This section records the nonlinear step between a local Gaussian reader and
the group-valued field that it is meant to read. All comparisons use one
coarse field, one gauge based at one site, one parent forecast, and one joint
Gaussian sample. A patch restricts which addresses are tested; it does not
change the forecast, covariance, or gauge from address to address. The
local Gaussian hypotheses used here are those of
Section~\ref{sec:actual-reader-local-gaussian-prices}. That section's
same-pose mean and covariance estimates remain inputs, as do the paid
regularity and window hypotheses specified below.

Let $x$ be the large coarse parameter, $L$ the fixed blocking side, and $w$
the fixed width multiplier of the fine read. Set
\[
 y=x+c_L,\qquad R=x^{-3/8},\qquad
 g=\beta^{-1/2},\qquad g'=(\beta')^{-1/2}.
\]
Here $c_L\ge0$ is the fixed scale shift. The coupling band is
$x/2\le\beta\le3x$, $x/2\le\beta'$ and $|\beta'-\beta|\le\gamma$,
for a fixed $\gamma>0$.
In particular $\sqrt{x/2}\le g^{-1}\le\sqrt{3x}$. Define
\[
 \operatorname{phase}(v)=v^{-1/2}\exp((\log\log v)^{3/2}),
 \qquad \tau=\operatorname{phase}(x),\qquad t=\operatorname{phase}(y).
\]
Set $B(v)=\lceil(\log v)^2\rceil$ and $W_f=w B(y)^2$. The recording parameter
$r(x)\in[x,x^2]$ has temperature $\tau$. Thus the logarithmic clearance
allowance from Section~\ref{sec:ex-lattice-actual-reader-small-window}
on this recording scale is
\[
 a_x=16\{2(B(r(x))+1)+1\}^5 n_{\rm plaq}\,\tau .
\]
Fix the coarse chart constant $C_{\rm core}$ and a common constant
$C_0\ge1$ such that the deterministic link response to a direction of
size at most $a_x$ is bounded by $C_0a_x$ on the funded tube. Let
$C_{\rm high}$ and $A_*$ be the chord and curvature constants for the
two-exponential chart estimate proved below. Let $n_G$ be the real
dimension of the gauge algebra and put $e=(A_*+n_G)R^2$. All these
constants are fixed before taking large $x$.
Let $\mathcal C_o(V)_q$ denote the nonlinear plaquette-curvature coordinate
of the coarse field in the fixed gauge based at $o$, indexed by coarse
plaquettes and color coordinates. The parent forecast is the fixed real row
\[
 \mathcal P_{o,V,g'}(u,a)=\frac1{g'}
   \sum_q K_{u,a;q}\,\mathcal C_o(V)_q.
\]
Let $D_{A,F,x}$ denote the linear response on the selected store $A$
and forest $F$. It inserts the active coordinates of the joint Gaussian
vector into the fine link field, with zeros off the active coordinates,
then takes the signed curl and reads its real color coordinates in the
fixed frame. For a joint Gaussian vector $Z$ we write
\[
 \mathcal G_{A,F,x}(Z)=D_{A,F,x}Z
\]
for the resulting zero-padded signed-curl response.
When an incidence index $j$ is used instead of $(u,a)$, it denotes that
same plaquette/color row. The same $V,o,g'$ and kernel $K$ are used for every tested row. In
particular, $\mathcal P$ is not recomputed from a patch-dependent field.

\subsection{The scalar scale and ticket estimates}

Fix a patch of side $L$ and a finite transition tube with its row and
chart bounds. For the positive cut rate $s$ and nonnegative strip
parameter $z$ of this tube, put
\[
 m_2=\binom42=6,\qquad M_0=\frac{m_2^2 z^3}{L^2},
 \qquad P_0=m_2M_0.
\]
The two diagonal installations and the transport row give the cubic
strip bound $M_0$, as in Section~\ref{sec:actual-reader-short-covariance}.
The factor $m_2$ in $P_0$ sums over two-form output labels. For fixed
$z$ choose $L$ large enough that $2n_GP_0\le1/4096$, which follows from
$P_0=m_2^3 z^3/L^2$. Define the row-tail and comparison bounds
\[
 h_x=\left\lfloor\frac{B(x)}{10000}\right\rfloor,\quad
 \Delta(x)=P_0e^{-s h_x},\quad
 E(x)=C_{\rm core}\bigl(\Delta(x)R+P_0R^2\bigr).
\]
The row and chart estimates to which these bounds are applied are
specified in Lemma~\ref{lem:readerB-pointing}.

Since $r(x)\le x^2$, the logarithmic factor in $a_x$ is bounded by
a fixed power of $1+\log x$. Also $\tau=x^{-1/2+o(1)}$,
$t\ge y^{-1/2}$, and eventually $t\le2x^{-7/16}$. The phase ratio
$\tau/t\le\sqrt{y/x}\le2$ follows for large $x$ from the monotonicity
of the exponential factor on $y\ge x$. For every fixed chart radius
$r_G>0$, these bounds and $R=x^{-3/8}$ give, eventually,
\begin{equation}\label{eq:readerB-clearance}
 4C_0 a_x\le R\le r_G,\qquad x^{-7/8}\le R^2,\qquad
 \Delta(x)\le R.
\end{equation}
For the last inequality use $h_x\ge(\log x)^2/10000-1$ and $s>0$.
In particular $E(x)\le C_{\rm core}(1+P_0)R^2$. Uniformly in the
coupling band,
\[
 |g/g'-1|=\frac{|\beta'-\beta|}
                {\sqrt\beta(\sqrt{\beta'}+\sqrt\beta)}
       \le\gamma/x\le x^{-1/8}.
\]
Since $E(x)/g=O(x^{-1/4})$, $g^{-1}=O(\sqrt{x})$ and
$\tau=x^{-1/2+o(1)}$, both terms of the following quotient are
$o(x^{-1/16})$, uniformly in the coupling band. Hence
\begin{equation}\label{eq:readerB-ticket-bill}
 \frac{E(x)+P_0\tau x^{-1/8}}{g}\le x^{-1/16}
\end{equation}
eventually. The vector price is also eventual. Indeed
$n_GP_0\tau\le t/4096$ by the choice of $L$, and
$n_GE(x)\le t/4096$ for large $x$ because $R^2=o(t)$.

Write $c_G$ for the color-scale constant in the curvature product estimate
below and $\delta_G$ for the radius of the local curvature lower bound.
Use $M=C_{\rm high}R$, and for positive $g$ and any room radius $W$ define
\begin{equation}\label{eq:readerB-transport-cost}
 \mathcal N_G(M,g,W)
 =\frac{c_G\{u(4M)+(u+4M)u+6M^2\}}{g},
 \qquad u=4WM.
\end{equation}
Section~\ref{subsec:readerB-gauge-read} derives this expression from the
based holonomy and the four-link product. For fixed $w$,
$W_f^2=O((1+B(x))^4)$ because $y=x+c_L$. Thus
$\mathcal N_G(M,g,W_f)=O((1+B(x))^4x^{-1/4})=o(x^{-1/8})$.
Together with $R^2=o(t)$ and $e/g=O(x^{-1/4})$, this gives the remaining
large-$x$ comparisons
\begin{align}
 4M&\le\delta_G,\label{eq:readerB-small-chart}\\
 6c_G M^2+4e&\le t/64,\label{eq:readerB-four-budget}\\
 \mathcal N_G(M,g,W_f)+4e/g&\le x^{-1/8}.\label{eq:readerB-nautical}
\end{align}

We use the local curvature estimate \eqref{eq:seed-curvature-bound}.
With the group-dependent radius $\delta_G$ and
$\operatorname{chord}(U)\le\delta_G$, it gives
\[
 \frac34\,\operatorname{chord}(U)^2
       \leq\|\kappa_G(U)\|^2.
\]
The constants in this inequality depend only on the fixed compact
gauge group. All large-$x$ thresholds here are for fixed $L$, patch,
and transition tube with the stated bounds.

\begin{lemma}[Eventual local room]
\label{lem:readerB-enlightened}
For fixed $G,L,w,z,s,\gamma$ as above and a transition tube with the stated
paid chart bounds, all sufficiently large $x$ and all
$x/2\leq\beta\leq3x$, $x/2\le\beta'$ and
$|\beta'-\beta|\le\gamma$ have the following properties. Set $\kappa_0=1/4096$.
\begin{enumerate}
\item the chart clearance is at most $R$, $R$ lies below the fixed chart
   margin, and $x^{-7/8}\leq R^2$;
\item for every coarse direction of norm at most the clearance, the sum of
   the deterministic response bound, $n_Gx^{-7/8}$ and the Gaussian
   reserve $\kappa_0t$ is at most $R$;
\item the near-identity link price is at most $C_{\rm high}R$, and the four-face
   curvature budget is at most $t/64$;
\item the sum of the four-face gauge-reading error, the coupling error and
   the residual curvature price is at most $x^{-1/32}$.
\end{enumerate}
The assertions hold uniformly in the period for the fixed patch and local
data.
\end{lemma}
\begin{proof}
The deterministic response has size at most $C_0a_x\le R/4$ by
\eqref{eq:readerB-clearance}.
Since
$x^{-7/8}/R=x^{-1/2}$ and
$t/R\leq2x^{-1/16}$, the color-count and reserve terms are each at most
$R/4$ after increasing $x_0$. This proves the room estimate. The link price
is the fixed chart constant times $R$; the four-face budget is
$O(R^2)=O(x^{-3/4})$, while $t\geq y^{-1/2}\geq c x^{-1/2}$ for large
$x$, proving the curvature budget. Finally, the definition of the
four-face gauge error is
\[
 \mathcal N_G(C_{\rm high}R,g,W_f)+4e/g.
\]
The scalar scale estimate bounds this by $x^{-1/8}$. Adding the paid row
read error $x^{-1/16}$ and the remaining coupling allowance gives at most
$x^{-1/32}$ eventually, since each stronger negative power divided by
$x^{-1/32}$ tends to zero. The fixed costs were chosen before the period,
so the threshold is period-independent.
\end{proof}

\subsection{From the coarse forecast to a fine curvature row}

Let the coarse side be $c$ and the fine side be $f$, with a fixed blocking
map. Let $V$ be the coarse field and $o$ a single base site. The common pose
assumption is that one based gauge transforms $V$ on the whole room about
$o$ into the exponential of one real direction $v$, with
$\|v\|_\infty\le a_x/2$, where the norm of a link field is the
maximum of its Euclidean color norms. The field is healthy at
threshold $\tau$ on that room. The curvature coordinate has norm at most the plaquette chord by
\eqref{eq:seed-normal-defect}, hence at most $\tau$ there. The chosen stores have no support there. The same gauge is used
for every row and plaquette below.

Let $F$ be the selected forest and $Z$ the actual joint Gaussian sample on
its active link coordinates. Write $z_j$ for its insertion into the fine
link field (zero outside the active coordinates). The stationary linear
fine response is $\ell=\mathcal U_Fv$ from Section~\ref{sec:local-mean},
using the same tube and forest.
On an internal bond set $b_j=0$. On a crossing bond set $b_j$ equal to the
corresponding coarse direction of $v$. The replica link is the identity
in the first case and copies the coarse link in the second. Because the
pose gives $V^o=\exp v$ at the relevant coarse bonds in the gauge based at
$o$, the replica link equals $\exp b_j$ in both cases, with
$\|b_j\|\le\|v\|_\infty\le R$. The stationary relative response of Section~\ref{sec:local-mean} is
$d_j=\ell_j-b_j$. This is the decomposition of the tangent in the
replica coordinates: differentiating $\exp(\epsilon d_j)\exp(\epsilon b_j)$
at $\epsilon=0$ gives $d_j+b_j=\ell_j$. The comparison mean
$m^0$ in \eqref{eq:reader-local-column-price} is the restriction of this $d$
to the active bonds in the same forest.

Let $w_j$ be the actual local mean in those relative coordinates and put
$r_j=gz_j$. The pointwise mean estimate
\eqref{eq:reader-local-column-price} at recording scale $r(x)\ge x$ bounds
each color component of $w_j-d_j$ by $x^{-7/8}$ on the interior active
bonds. The Euclidean norm is bounded by the sum of component moduli, so
$\|w_j-d_j\|\le n_Gx^{-7/8}$. The cleared stencils place every queried nonforest bond in the active set.
The update first uses the link increment $\exp(w_j+r_j)$ and then the
replica link. On forest bonds the update and the replica are both the
identity, so the relative coordinates $w_j,d_j,r_j$ there are zero. Thus the
actual fine field on each of these bonds is
\begin{equation}\label{eq:readerB-fine-link}
 U_j=\exp(w_j+r_j)\exp(b_j),\qquad \ell_j=d_j+b_j.
\end{equation}

Write $\mathsf M$ for the fixed unitary representation and
$\Pi_{\mathfrak g}$ for the real-linear color projection. In the physical
normalization its operator norm for the Frobenius norm is at most $1/c_G$,
while $\|\mathsf M(U)-I\|_F=c_G\operatorname{chord}(U)$.
The curvature coordinate is
\[
 \kappa_G(U)=\Pi_{\mathfrak g}\!\left(
       \frac{\mathsf M(U)-\mathsf M(U^{-1})}{2}\right)
             =\Pi_{\mathfrak g}(\mathsf M(U)-I).
\]
Here $\mathsf M(U^{-1})=\mathsf M(U)^*$, so removing the Hermitian part
does not affect the projection onto the skew-Hermitian Lie algebra. It
follows by expanding the product of the matrix increments that
\begin{equation}\label{eq:readerB-product-defect}
 \kappa_G(UV)-\kappa_G(U)-\kappa_G(V)
   =\Pi_{\mathfrak g}((\mathsf M(U)-I)(\mathsf M(V)-I)).
\end{equation}
The Frobenius product inequality bounds its norm by
$c_G\operatorname{chord}(U)\operatorname{chord}(V)$.

Choose $r_G>0$ so that the exponential and its curvature coordinate
have bounded second derivative on the closed ball of radius $r_G$.
The matrix exponential is smooth, and its derivative at zero is the
inclusion of the Lie algebra in the matrix space. Projecting this
inclusion back to color space is the identity. Taylor's formula
on radial segments therefore gives constants $C_{\exp},A_{\exp}>0$
depending only on the gauge chart such that
\[
 \operatorname{chord}(\exp a)\le C_{\exp}\|a\|,
 \qquad \|\kappa_G(\exp a)-a\|\le A_{\exp}\|a\|^2
 \quad (\|a\|\le r_G).
\]
For $\|a\|,\|b\|\le R\le r_G$, chord subadditivity bounds the chord of
$\exp a\exp b$ by $2C_{\exp}R$. Equation~\eqref{eq:readerB-product-defect}
and the two radial estimates bound the curvature error from $a+b$ by
$(c_G C_{\exp}^2+2A_{\exp})R^2$. Thus one may take
$C_{\rm high}=2C_{\exp}$ and $A_*=c_G C_{\exp}^2+2A_{\exp}$ in the
scalar estimates. Any increase of these fixed constants preserves
the same eventual powers.

The guarded event is
\[
 \|g z_j\|\leq\kappa_0t
\]
for every bond in the enlarged room. The paid room estimate is
\[
 \|w_j+r_j\|\leq \|d_j\|+n_Gx^{-7/8}+\kappa_0t\leq R,
\]
and we have shown $\|b_j\|\le R$. Apply the two-exponential chart
bound to $(w_j+r_j,b_j)$ to obtain
\begin{align}
 \operatorname{chord}(U_j)&\leq C_{\rm high}R,\label{eq:readerB-bondchord}\\
 \|\kappa_G(U_j)-(\ell_j+r_j)\|&\leq e.
 \label{eq:readerB-bondcurv}
\end{align}
Indeed, $w_j+r_j+b_j=(d_j+b_j)+r_j+(w_j-d_j)=\ell_j+r_j+(w_j-d_j)$.
The portrait term is at most $n_Gx^{-7/8}$; since both Taylor arguments
have norm at most $R$, the two-link curvature remainder is at most $A_*R^2$. Since $x^{-7/8}\leq R^2$, their sum is
bounded by $e=(A_*+n_G)R^2$. Here $\ell_j+gz_j$ is the value of the link direction at the single bond
$j$; the face boundary sum is taken only in the next subsection.

\begin{lemma}[Read error against the common parent forecast]
\label{lem:readerB-pointing}
Assume the paid row and transport estimates for the selected transition
tube: on every readable address in the room, the fine response row differs
from its shorter parent row by total absolute mass at most $\Delta$, the
parent row has mass at most $P_0$, and the deterministic coarse curvature
Taylor bounds stated below hold. Let
$\Delta=m_2 M_0 e^{-s h_x}$, the absolute tail bound for the
two-form row with strip bound $M_0$. Then for every interior fine
plaquette $u$ and color coordinate $a$, with
$f=\operatorname{faceCurl}(u,\ell)$,
\begin{equation}\label{eq:readerB-pointing-conclusion}
 \|f\|\leq t/2048,\qquad
 \left|g^{-1}f_a-\mathcal P_{o,V,g'}(u,a)\right|
      \leq x^{-1/16}.
\end{equation}
\end{lemma}
\begin{proof}
Index coarse plaquette/color coordinates by $q$. Let $m_q$ be the
linearized coarse plaquette curvature in direction $v$, and let
$(u_0)_q=\mathcal C_o(V)_q$ be the nonlinear curvature in the same based
gauge. The exact response identities are
\[
 f_a=\sum_q P_{u,a;q}m_q,
 \qquad Q_{u,a}=\sum_qK_{u,a;q}(u_0)_q,
 \qquad \mathcal P_{o,V,g'}(u,a)=Q_{u,a}/g'.
\]
The deterministic chart estimates give the global linear-row bound
$|m_q|\leq C_{\rm core}R$ for every tested coarse plaquette/color
coordinate. On the support of the parent row $K_{u,a;\cdot}$ they also give
$|m_q-(u_0)_q|\leq C_{\rm core}R^2$.
The purchased row bounds are
\[
 \sum_q|P_{u,a;q}-K_{u,a;q}|\leq\Delta,
 \qquad \sum_q|K_{u,a;q}|\leq P_0.
\]
Subtract the response sums and add/subtract $\sum_qK_{u,a;q}m_q$. Then
\[
 |f_a-Q_{u,a}|
 \leq C_{\rm core}R\sum_q|P_{u,a;q}-K_{u,a;q}|
       +C_{\rm core}R^2\sum_q|K_{u,a;q}|
 \leq \Delta C_{\rm core}R+P_0C_{\rm core}R^2
 =E(x).
\]
For the normalized comparison, $|Q_{u,a}|\leq P_0\tau$ since the
nonlinear parent curvature coordinates have modulus at most $\tau$. Thus
\[
 \left|\frac{f_a}{g}-\frac{Q_{u,a}}{g'}\right|
 \leq\frac{|f_a-Q_{u,a}|+|Q_{u,a}|\,|g'^{-1}g-1|}{g}
 \leq\frac{E(x)+P_0\tau x^{-1/8}}g
 \leq x^{-1/16},
\]
by the coupling bound and \eqref{eq:readerB-ticket-bill}. For the vector
bound, sum the component bounds to get
$\|f\|\le n_G(P_0\tau+E(x))\le t/2048$ by the scalar estimates above.
All rows have
finite support in the selected window, and the bounds are period-uniform.
\end{proof}

\subsection{The four-bond curvature and the common gauge read}

For a fine plaquette $u$, order its bonds $j_1,j_2,j_3,j_4$ so that the
plaquette product is
\[
 \Xi_u=U_{j_1}U_{j_2}U_{j_3}^{-1}U_{j_4}^{-1}.
\]
For a link field $A$, let $\operatorname{faceCurl}_u(A)$ be the signed
sum of the four link coordinates in this order.
Apply the two-factor relation above successively to the four oriented
factors $U_{j_1},U_{j_2},U_{j_3}^{-1},U_{j_4}^{-1}$. Each has chord at most
$M=C_{\rm high}R$, and inversion preserves chord while negating curvature.
At the second factor the defect is at most $c_GM^2$. At the third, the
preceding product has chord at most $2M$, so the new defect is at most
$2c_GM^2$. At the fourth, its chord is at most $3M$, so the new defect is
at most $3c_GM^2$. Summing gives exactly
$(1+2+3)c_GM^2=6c_GM^2$. Replacing the four
individual link-curvature coordinates by their linear directions using
\eqref{eq:readerB-bondcurv} costs at most $4e$. Consequently
\begin{equation}\label{eq:readerB-four-product}
 \|\kappa_G(\Xi_u)-\operatorname{faceCurl}_u(\ell+gz)\|
       \leq 6c_G(C_{\rm high}R)^2+4e\leq t/64,
\end{equation}
by \eqref{eq:readerB-four-budget}. Linearity of the face curl, the bound
from Lemma~\ref{lem:readerB-pointing}, and the four guarded link increments
give
\[
 \|\operatorname{faceCurl}_u(\ell+gz)\|
       \leq t/2048+4\kappa_0t.
\]
Consequently the product curvature coordinate has norm at most
$t/2048+4\kappa_0t+t/64<t/32$.
Since $\kappa_0=1/4096$,
\[
 \frac1{2048}+\frac4{4096}+\frac1{64}<\frac1{32}.
\]
The plaquette chord is at most four times the link chord, so
$4C_{\rm high}R\leq\delta_G$ keeps it inside the local curvature chart.
The local lower bound $\tfrac34\operatorname{chord}(\Xi_u)^2
\le\|\kappa_G(\Xi_u)\|^2$ therefore gives
\begin{equation}\label{eq:readerB-healthy-face}
 \operatorname{chord}(\Xi_u)\le\frac{2t}{32\sqrt3}<t/8.
\end{equation}
Thus if the health flag threshold is the phase $t$, the tested flag is
strictly below threshold. Applying this to each plaquette in the tested
patch proves the healthy-mask conclusion. It is a local conclusion on the
guarded sample; all plaquettes used in this deduction lie in the enlarged room.

\subsection{Based transport for the fine read}
\label{subsec:readerB-gauge-read}

Let $r$ be the read center on the fine torus and $v_0$ the base vertex
of the tested plaquette $u$. The read is in the based gauge obtained
by transporting along the fixed coordinate paths from $r$. If the fine
period is $n$, the condition $2(W_f+1)<n$ gives an unwrapped
coordinate box of radius $W_f$ about $r$. When the plaquette is at least
two steps inside this box, its four bonds and the coordinate path
to $v_0$ lie in it. The latter path has at most $4W_f$ steps since
each of the four displacements is at most $W_f$. Write $H$ for its
ordered product in $U$. By the chord triangle and inverse identities,
$\operatorname{chord}(H)\le 4W_fM$ and
$\operatorname{chord}(\Xi_u)\le4M$.

For a site gauge $k$ sending $U'$ to $U$, multiplying the gauge-transformed
links along the path cancels the gauge factors at all intermediate sites.
The path product transforms by $k(r)$ at the start and $k(v_0)^{-1}$ at
the end. At the plaquette the factors telescope to conjugation by
$k(v_0)$. Hence the based plaquette in $U$ is the based plaquette
in $U'$ conjugated by $k(r)$. The curvature projection is equivariant
under conjugation, since the adjoint action of the unitary representation
preserves the Lie subspace and its orthogonal complement. It follows that
the normalized curvature read in $U'$ after the color rotation by $k(r)$
equals $g^{-1}\kappa_G(H\Xi_uH^{-1})$ in $U$.

Apply \eqref{eq:readerB-product-defect} at $(H,\Xi_u)$ and
$(H\Xi_u,H^{-1})$. The terms $\kappa_G(H)$ and $\kappa_G(H^{-1})$
cancel. With $u_0=4W_fM$, the two defects give
\[
 \|\kappa_G(H\Xi_uH^{-1})-\kappa_G(\Xi_u)\|
 \le c_G\{u_0(4M)+(u_0+4M)u_0\}.
\]
Before division by $g$, the difference between $\kappa_G(\Xi_u)$
and the signed sum of the four bond curvature coordinates is bounded
by $6c_GM^2$, as proved above. This yields exactly
$\mathcal N_G(M,g,W_f)$ in \eqref{eq:readerB-transport-cost} as an upper bound
for the discrepancy between the rotated read and the normalized signed sum.
Replacing each bond curvature by $\ell_j+gz_j$ costs at most $e/g$
per bond by \eqref{eq:readerB-bondcurv}. The signed sum then reads
$g^{-1}f_a+\mathcal G_{A,F,x}(Z)(u,a)$ in color coordinate $a$, by linearity
and the definition of $D_{A,F,x}$. Finally
\eqref{eq:readerB-pointing-conclusion} replaces $g^{-1}f_a$ by the common
parent forecast with error $x^{-1/16}$. The total is at most
\[
 \mathcal N_G(M,g,W_f)+4e/g+x^{-1/16}
 \le x^{-1/8}+x^{-1/16}\le x^{-1/32}
\]
for large $x$. The gauge used to compare $U'$ to $U$ can vary. The
within-room forecast, sample, and coarse pose are fixed throughout.

\begin{proposition}[Actual-reader curvature and read estimate]
\label{prop:readerB-curvature-read}
Fix a selected transition tube and one posed coarse field with based gauge
at $o$. Suppose the actual Gaussian reader estimates of
Section~\ref{sec:actual-reader-local-gaussian-prices}, the paid row and chart
bounds above, and the local window/separation conditions hold, including
$2(W_f+1)<n$ for the fine period $n$. For each sample in the
enlarged room with $\|gz_j\|\le\kappa_0t$ on every active bond,
the resulting fine field is healthy on the selected patch. For every gauge-equivalent fine field and every incidence row $j=(u,a)$
whose two-step neighborhood lies inside the patch, its rigid color read differs from the same parent
forecast plus the transported Gaussian curvature by at most
$x^{-1/32}$. All conclusions are uniform in the period once the displayed
fixed-data threshold is reached.
\end{proposition}
\begin{proof}
Use \eqref{eq:readerB-fine-link} for all bonds in the enlarged room,
then apply Lemma~\ref{lem:readerB-pointing} to each tested row.
Equation~\eqref{eq:readerB-healthy-face} gives the health estimate.
The based-transport calculation and \eqref{eq:readerB-nautical} give
the read estimate. All statements use
the same base site and parent forecast, and their constants do not depend
on the number of sites outside the tested patch.
\end{proof}

\subsection{The lost-mask estimate for a correlated Gaussian}

Let $\mathcal I_y$ be the actual finite set of real link-color coordinates
read by the guard on the enlarged room. Its cardinality satisfies
\begin{equation}\label{eq:readerB-room-dimension}
 m_y:=|\mathcal I_y|\leq C_w(\log y)^{16}
\end{equation}
for a constant $C_w$ depending on the fixed width and color dimension, but
not on the period. Indeed the room radius for the mask is $W_f+8$,
$B(y)^2\le4(\log y)^4$ for large $y$, and the number of
link-color coordinates in a four-dimensional box is at most a fixed
multiple of its side to the fourth power.

If $m_y=0$, the product veil is $1$ and the loss is zero. Otherwise assume
$m_y\geq1$; the gauge algebra is nontrivial, so $n_G\geq1$. Group the
coordinates by bond, writing $X_b\in\mathfrak g$ for the actual Lie-valued
fluctuation read on bond $b$. Let $X=(X_i)_{i\in\mathcal I_y}$
have the actual centered joint Gaussian law. Its covariance matrix is
positive semidefinite and each entry has
absolute value at most $U$, where $U\geq1$ is the common palette bound from
the actual reader estimates. For $v>0$, each marginal obeys the Chernoff
bound
\[
 \mathbb P(|X_i|\geq v)\leq2\exp(-v^2/(2U));
\]
this follows by optimizing
$\mathbb E e^{sX_i}=e^{s^2\operatorname{Var}(X_i)/2}$ at $s=v/U$ and
applying the same estimate to $-X_i$. A union bound, which does not assume
independence, gives
\[
 \mathbb P\{\max_i|X_i|\geq v\}
       \leq2m_y\exp(-v^2/(2U)).
\]
For the bondwise guard, consider the event $\max_b\|X_b\|\ge v$.
On this event one of the $n_G$ real coordinates of some $X_b$
has modulus at least $v/\sqrt{n_G}$. A union bound therefore
gives
\[
 \mathbb P\{\max_b\|X_b\|\geq v\}
       \leq2m_y\exp(-v^2/(2n_GU)).
\]
Set $\varrho_y=[8n_G(1+m_yU)]^{-1}$. Then
$1/(2n_GU)-\varrho_y\geq3/(8n_GU)$. The guard threshold grows faster than
every power of $\log y$, so eventually
$3v^2/(8n_GU)\geq\log(2m_y)-2$. Hence
\begin{equation}\label{eq:readerB-correlated-reserve}
 \mathbb P\{\max_b\|X_b\|\geq v\}
       \leq e^2\exp(-\varrho_y v^2).
\end{equation}
The vector is the actual joint Gaussian. Correlation is allowed throughout;
only its marginal variance bound and the finite union bound are used.

Take the guard width
\[
 v=\frac{\kappa_0 t}{2g}.
\]
For the band $x/2\leq\beta\leq3x$, the scalar coupling bound gives
$v\geq\delta_0\,\operatorname{hot}(y)$ with a fixed $\delta_0>0$;
here $\operatorname{hot}(y)=\exp((\log\log y)^{3/2})$. By
\eqref{eq:readerB-room-dimension},
\[
 1+m_yU\leq (1+C_wU)(\log y)^{16}.
\]
For every fixed $p\geq0$, the ratio
\[
 \frac{\delta_0^2\exp(2(\log\log y)^{3/2})}
      {8n_G(1+C_wU)(\log y)^{16}(2+p)\log y}
\longrightarrow\infty.
\]
Indeed, writing $q=\log\log y$, the logarithm of the numerator growth
relative to the denominator is $2q^{3/2}-17q+O(1)$, which tends to
$+\infty$. Thus eventually
$\varrho_y v^2\geq(2+p)\log y$, and
\eqref{eq:readerB-correlated-reserve} gives
\begin{equation}\label{eq:readerB-tail-power}
 \mathbb P\{\max_b\|X_b\|\geq v\}
       \leq e^2 y^{-(2+p)}\leq x^{-p}.
\end{equation}

Let $\Gamma$ be the smooth cutoff on the actual Gaussian vector: one
radial factor for each active bond in the enlarged room. Each factor
lies in $[0,1]$, vanishes outside the vector guard radius $\kappa_0t/g$,
and equals one on the ball of half that radius. Consequently
\[
 0\leq1-\Gamma(X)\leq
      \mathbf1_{\{\max_b\|X_b\|\geq v\}}.
\]
Integrating and using \eqref{eq:readerB-tail-power} proves
\[
 \mathbb E[1-\Gamma(X)]\leq x^{-p}.
\]
The estimate uses the marginal variances of the actual correlated
Gaussian on its finite room.

\begin{lemma}[One-patch Gaussian guard]
\label{lem:readerB-mask}
For every fixed $p\geq0$, fixed patch width, and common palette entry bound
$U\geq1$, the mask loss on the actual joint Gaussian sample is at most
$x^{-p}$ for all sufficiently large $x$, uniformly in the period. In
particular, choose $p=1$ to obtain the mask loss $x^{-1}$ used in the
reader purchase.
\end{lemma}
\begin{proof}
The coordinate count is \eqref{eq:readerB-room-dimension}; the one-coordinate
Chernoff estimate, union bound and scale comparison establish
\eqref{eq:readerB-tail-power}. The pointwise cutoff comparison then proves
the expected mask loss.
\end{proof}

\begin{remark}[Boundary of this package]
The result proves the nonlinear curvature and read estimates and the power-small
loss of the Gaussian cutoff for the fixed posed reader, under the stated
Gaussian, chart, row and window hypotheses. The assertion that the
warning factor vanishes for every compatible field/store palette requires
the separate hypothesis of Section~\ref{sec:actual-corrected-body-compiler}.
\end{remark}
